\documentclass[10pt, a4paper]{memoir}

\usepackage[utf8]{inputenc}
\usepackage[T1]{fontenc}
\usepackage[polish]{babel}

\usepackage{graphicx}
\usepackage{caption}
\usepackage{verse}
\usepackage{amsmath}

\setlength{\vleftmargin}{0pt}
\indentpattern{0}
\pagestyle{plain}
\linespread{1.3}

\setsecnumdepth{none}
\setlength{\leftmargini}{0pt}

\title{.gitignore}
\author{Kate Malec}
\date{\the\year}

\begin{document}

\maketitle
\thispagestyle{empty}
\newpage

\thispagestyle{empty}
\vspace*{\fill}
\begin{flushright}
    \itshape
    To wszystko to tylko literacka fikcja
\end{flushright}

\newpage
\tableofcontents
\newpage

\section{Można być inteligentnym, a zarazem chamem i~prostakiem}

Mój ojciec przychodzi po dzieci. Wchodzi do domu. Jest ostatnio grzeczny, odkąd sprzedałam mu parę solidnych pstryczków w~nos. W~końcu też chyba zaczął się starzeć i~zaczyna mieć radość ze wspólnego czasu z~wnukami. Nie wiem, coraz mniej mnie to już interesuje. To jemu kończy się czas~-- nie mnie.

Wchodzi, od drzwi czuję jego ego, które jednak w~miarę możliwości on stara się trzymać na wodzy. Mimo wszystko wparowuje do domu tak, jak stoi~-- w~butach. Nie zwracam mu już uwagi, uznałam, że trudno. Myślę jednak o~tej jednej znajomej, która tak pilnowała, aby w~odwiedzinach ściągać buty przed progiem. Wzruszam ramionami, nie mam już siły na takie pierdoły.

Ojciec zamaszyście przechodzi przez mieszkanie. Wyczuł chyba jednak moje pole, bo pyta: „czy można tak w~butach”. Mówię, że niekoniecznie. Ściąga więc. Po drodze mija stertę obrazów już na wejściu, dalej kilka na szafce. Nie zauważa jednak nic. Coś gada, przeważnie o~pracy i~o~tym, jak obronną ręką wyszedł z~trudnej sytuacji. Jakiej trudnej sytuacji? Pracując w~dwóch zagranicznych korpach za taką kasę, jaką ludzie oglądają tylko na filmach?

Dzieci jedzą. Namyśliły się na żarcie tuż przed przyjazdem dziadka. To im zawdzięczam jego wizytę w~domu. Nie miałabym nic przeciwko, w~końcu to mój ojciec. Nie mogę się jednak pozbyć tego napięcia. Zawsze, kiedy się widzimy, moje pole elektryczne jeszcze przez kilka godzin będzie najeżone. Będę się złościć o~pierdoły i~mieć reakcje kompulsywne. Nigdy chyba z~tego nie wyrosnę, choć trwa to już teraz znacznie krócej i~ma łagodniejszy przebieg. Namiętnie pogrążam się w~zmywaniu naczyń, a~większość dialogu odbywam odwrócona do niego plecami.

Ojciec, mówiąc wciąż o~sobie, przechadza się przez mieszkanie. Zatrzymuje się przed obrazem~-- grafiką kazuara, który wisi tutaj już od ponad roku. Widział go wielokrotnie, nigdy nie zwracając uwagi. Podchodzi i~zamaszystym tonem podszytym dozą kpiny i~ignorancji czyta podpis, myląc się przy tym:

-- Kazaur chełmiasty.

-- Kazuar -- poprawiam go znad~zlewu. -- Ale nie jestem pewna, może to ja popełniłam literówkę.

Mamrocze pod nosem. Wskazuje stojące w~kącie deski, które mąż kupił z~myślą o~zrobieniu półek:

-- No, to teraz będzie poważniejsze malowanie!

Uśmiecham się tylko pod nosem. Myślę sobie o~teorii dusz, z~którą się ostatnio poznałam. Podobno sama wybrałam sobie tego człowieka na ojca. Tak jak matkę. Myślę, że bardzo kuszący musiał być dla mnie wykres natalny, który mi sprezentowali. Trzymam się tej teorii, bo nie znajduję innego wytłumaczenia na nasze powiązanie. Jak w~ogóle moja matka z~nim wytrzymuje? Upominam się: nie mogę być dla niego zbyt surowa. To dzięki niemu mam rozbudowaną funkcję kognitywną \textit{Ti}. Błogosławione \textit{T} w~magicznym czteroliterowym skrócie. Dzięki niej potrafię tak beznamiętnie ocenić jego błazeństwo i~konfrontować się z~nim, przełączając w~tryb bezemocjonalnej suki. Zastanawiam się, czy gdyby wychowywał mnie ktoś inny, też bym tak umiała? Czy jest to cecha, którą wykształciłam w~dzieciństwie? Nigdy się tego nie dowiem na pewno.

Myślę sobie, że ta marża na hasło \textit{rodzina}, którą postanowiłam doliczać wszystkim spokrewnionym, jest nieadekwatna. Zamiast 1.5k powinnam im naliczać 1.5\%.

\newpage

\section{Fuck around and find out}

Obiecałam sobie, że nie zdeprawuję żadnego pracownika uczelni.

Nie byłoby to wcale takie trudne, bo moje libido już od dłuższego czasu oscyluje wokół zera. Na dodatek, co było dla mnie ogromnym zaskoczeniem, wydział artystyczny obfituje w~nauczycieli różnej rangi, ale wszyscy przeważnie mają jedną wspólną cechę: rozbudowaną funkcję \textit{Fi}. A~to niestety jest dla mnie dyskwalifikujące~-- niezależnie od stopnia naukowego, jaki się uprawia, nie mogłabym się oprzeć, by odrobinę nie zdominować tego nieszczęsnego \textit{F}.

Mój mąż mnie za to nienawidzi. Kiedy tylko mówię, że właściwie wszyscy nauczyciele kierują się uczuciami jako głównym wyznacznikiem, patrzy na mnie ostro. Niemal z~pogardą. Ale co ja mogę na to poradzić? Ogólnie po studiach ścisłych ten wydział artystyczny to jak herbata i~ciastko.

Jest jednak wśród pracowników jeden nurt, który wyłamuje się mojemu zawyżonemu poczuciu wartości: pracownicy techniczni pracowni. Jest taki jeden, z~którym mam niemal natychmiastowe zdalne połączenie. \textit{ISTP}~-- nie mam co do tego najmniejszych wątpliwości. Dużo iestepeków chodzi po tym świecie w~spodniach roboczych. Niestety niewielu z~nich udźwignęłoby moją ciemną, intelektualną stronę. 

Obiecałam sobie jednak, że choćby nie wiem co, nie będę siać rozpusty na~uczelni. Później będzie, że osiągnęłam to, co osiągnęłam, przez łóżko. I~jeszcze ktoś na mojej pracy zyska za free. Może traktowałam to powiedzenie nieco zbyt dosłownie?
\newpage
\section{Matki i czarownice}
Dziewczyny z grupy zwracają się do mnie per \textit{matko}. Nie byłoby może w tym nic dziwnego, gdyby nie fakt, że dzieli nas zaledwie dekada. Mimo wszystko lgną do mnie i pomimo że sama ledwo wyrabiam na zakrętach, nie mam sumienia im odmówić. Lgną do mnie także i te starsze, ale dla nich jestem chyba nieco zbyt szybka, a na pewno --- one dla mnie za wolne. Najbardziej pasuje mi towarzystwo tegorocznych maturzystek. 

To grono jest dla mnie wyjątkowe, choć składa się na nie także kilka osób, które pierwsze przygody ze studiami mają już za sobą. Nawet się tego nie spodziewając, znalazłam się w środowisku, gdzie moje podejście \textit{hop do przodu} nikomu nie zawadza. Wprost przeciwnie. Dziś usłyszałam, że jestem najlepszym sołtysem. Jest to naprawdę zupełnie nowe doświadczenie.

Na studiach jest też Ula. Z Ulą znamy się z podyplomowych, gdzie ona, podobnie jak ja, została gorąco zaproszona na studia dzienne. Obydwie mamy spore ambicje oraz po garstce dzieci, choć jej są starsze. Nie wiem, ile dokładnie ma lat, ale życzę sobie wyglądać w jej wieku tak dobrze, jak ona. 

Oprócz określenia \textit{matka}, z Ulą zostałyśmy ochrzczone jeszcze jednym terminem: \textit{witch}. Przylgnął do nas zupełnie niewinnie, kiedy, rozmawiając o zdrowiu, koleżanki zaczęły się otwierać na swoje problemy, na co ja przyniosłam mieszankę ziół i powiedziałam, gdzie mogą się zaopatrzyć w remedia. Nigdy nie byłam zbyt mocna w zielarstwie, ale przerobiłam na tyle długą listę różnych dolegliwości, że z głowy wyrecytowałam solidną listę rozwiązań. Ula dopowiedziała coś od siebie i tak się zaczęło.

Z Ulą siedzimy w jednej ławce na wykładzie z technologii malarstwa. Prowadzący jest młody, choć ma już tytuł doktora. Gdy mu się przyjrzałam, podejrzewam, że może być młodszy ode mnie. Patrzę, jak ten wysoki, a przygarbiony już gość, chodzi i mamrocze pod nosem wątki poboczne wykładu kierowanego do sali pełnej ludzi, i czuję, że powinien mi się podobać. Myślę, jak bardzo przypomina mi uroczego Foxa Muldera i postanawiam, że tak będę go nazywać. Jest pasjonatą swojego przedmiotu, a przy tym --- totalnie odalienowany i fatalnie się komunikuje. Wykapany Fox. Słuchając go, trącam łokciem Ulę i szepczę w jej kierunku:

--- Nie może być żonaty.

Ula natychmiast przytakuje.

--- Nie może. 

Obydwie się śmiejemy.

Mówię do Uli, że jest słodki i wzbudza we mnie instynkt macierzyński. Ula, śmiejąc się, przytakuje. Jednocześnie zauważamy, że zupełnie naturalnie prowadzącego popchnęło w stronę jednej ze studentek. Zwraca się często wprost do niej, zaczyna wyraźniej mówić. Trącamy się z siostrą --- już jest sprzedany, możemy o nim zapomnieć. W zasadzie nie jest mi nawet przykro, a właściwie to się cieszę. Przekazuję dziewczynom z tyłu, że sorry, ale prowadzący jest już zajęty i wyjaśniam sytuację. Wzdychają. Jedna zauważa, że to nawet dobrze, bo ona woli tego z innej pracowni. Pytam --- którego? No i okazuje się, że tego, którego ja sama miałam na oku. Czuję jednak, że jej się należy bardziej. Poza tym mnie to już chyba w życiu wystarczy. Rozmowa kończy się tym, że pokazuję im jeszcze zdjęcia mojego męża. Zastanawiam się tylko, czy na pewno dobrze robię? Przecież nie życzę im źle.

Tak oto staję się czarownicą znacznie szybciej, niż zakładałam.
\end{document}














